\chapter{Biographical Sketches}
\label{app:biographies}

This appendix provides short biographies of the major figures discussed in this book.

\section*{Henri Poincar\'{e} (1854--1912)}

Jules Henri Poincar\'{e} was a French mathematician, theoretical physicist, engineer, and philosopher of science. Often described as a polymath and the last universalist in mathematics, he made fundamental contributions to pure and applied mathematics, including celestial mechanics, topology, number theory, and the theory of functions. He introduced the concept of the fundamental group and formulated the Poincar\'{e} conjecture in 1904. He also independently discovered the special theory of relativity before Einstein.

\section*{Grigori Perelman (1966--)}

Russian mathematician who proved the Poincar\'{e} conjecture and the Thurston geometrization conjecture. Born in Leningrad, he showed exceptional mathematical talent from an early age, winning a gold medal at the International Mathematical Olympiad with a perfect score. After completing his PhD at the Steklov Institute, he held positions in the United States before returning to Russia. In 2002--2003, he posted his proof online. He declined both the Fields Medal (2006) and the Clay Millennium Prize (2010). He has since withdrawn from public life.

\section*{Bernhard Riemann (1826--1866)}

German mathematician whose work revolutionized geometry, analysis, and number theory. His 1851 doctoral dissertation founded Riemannian geometry; his 1859 paper on prime numbers introduced the zeta function and the Riemann hypothesis. Despite his short life (he died of tuberculosis at 39), his ideas shaped much of modern mathematics and physics---Einstein's general relativity uses Riemannian geometry.

\section*{Stephen Cook (1939--)}

American computer scientist and mathematician. In his 1971 paper ``The complexity of theorem-proving procedures,'' he introduced the concept of NP-completeness and proved the Cook--Levin theorem, establishing SAT as the first NP-complete problem. He won the Turing Award in 1982.

\section*{Leonid Levin (1948--)}

Soviet-American mathematician and computer scientist. Independently of Cook, he developed the theory of NP-completeness in the Soviet Union (1973). His work on average-case complexity and algorithmic information theory has been highly influential.

\section*{Claude-Louis Navier (1785--1836)}

French engineer and physicist. He derived the Navier--Stokes equations in 1822, working from both an engineering and a mathematical perspective. He also contributed to the theory of elasticity.

\section*{Sir George Gabriel Stokes (1819--1903)}

Irish mathematician and physicist. He placed the Navier--Stokes equations on a rigorous physical foundation. His mathematical work spanned fluid dynamics, optics, and mathematical physics. He served as President of the Royal Society and as MP for Cambridge University.

\section*{Chen-Ning Yang (1922--)}

Chinese theoretical physicist. Together with Robert Mills, he developed Yang--Mills theory in 1954. He won the Nobel Prize in Physics in 1957 (with Tsung-Dao Lee) for work on parity violation. His contributions to statistical mechanics and gauge theory are foundational.

\section*{Robert Mills (1927--1999)}

American physicist who, as a postdoctoral researcher at Brookhaven National Laboratory, collaborated with Yang on the theory that bears their name. He spent most of his career at Ohio State University.

\section*{William Vallance Douglas Hodge (1903--1975)}

Scottish mathematician who developed the theory of harmonic integrals and the Hodge decomposition. He held the Lowndean Chair of Astronomy and Geometry at Cambridge and made fundamental contributions to algebraic geometry and differential geometry.

\section*{Pierre Deligne (1944--)}

Belgian mathematician, Fields Medalist (1978). He proved the Weil conjectures, revolutionizing algebraic geometry. He wrote the official description of the Hodge conjecture for the Clay Institute and made deep contributions to Hodge theory.

\section*{Bryan Birch (1931--)}

British mathematician known for the Birch and Swinnerton-Dyer conjecture. His work on elliptic curves and modular forms has been highly influential. He was a student of J. W. S. Cassels.

\section*{Sir Peter Swinnerton-Dyer (1927--2018)}

British mathematician and academic administrator. With Birch, he formulated the BSD conjecture based on numerical experiments on the EDSAC computer. He served as Master of St Catharine's College, Cambridge, and chaired the Higher Education Funding Council for England.

\section*{Andrew Wiles (1953--)}

British mathematician who proved Fermat's Last Theorem in 1994, one of the greatest achievements in number theory. He wrote the official description of the Birch and Swinnerton-Dyer conjecture for the Clay Institute. He was knighted in 2000.

\section*{Charles Fefferman (1949--)}

American mathematician. He became the youngest full professor in US history at age 22 and won the Fields Medal in 1978 for his work in analysis. He wrote the official description of the Navier--Stokes problem for the Clay Institute.

\section*{Arthur Jaffe (1937--)}

American mathematical physicist. He made fundamental contributions to constructive quantum field theory and served as the first president of the Clay Mathematics Institute. He co-wrote the Yang--Mills problem description with Edward Witten.

\section*{Edward Witten (1951--)}

American theoretical physicist and Fields Medalist (1990). He is widely regarded as the leading figure in string theory. His work has connected quantum field theory to deep areas of mathematics, including knot theory and mirror symmetry.

\section*{John Milnor (1931--)}

American mathematician who won the Fields Medal (1962), the Wolf Prize (1989), and the Abel Prize (2011). He wrote the official description of the Poincar\'{e} conjecture for the Clay Institute. His work spans topology, geometry, and dynamical systems.
